\begin{exercisechunk}
\item \textbf{Parallel scans of affine recurrences.}
\label[exercise]{ex:l5-affine-scan}
\exerciseprofile{3}{2}{2}{\exproof\quad\excalculation\quad\exinterpretation}{%
Derive the work, depth, and storage costs of affine scans and identify the
advantage of diagonal transitions.}
For affine maps, define
\[
 (A_2,b_2)\circ(A_1,b_1)=(A_2A_1,A_2b_1+b_2).
\]
Verify that this operation is associative. Show that a dense
$d$-dimensional sequential recurrence has $O(Td^2)$ work and $O(T)$
sequential depth, while a work-inefficient dense parallel scan has
$O(Td^3)$ work and $O(\log T)$ composition depth. Compare their activation
storage. Then show that diagonal transitions reduce the scan to $O(Td)$ work,
$O(\log T)$ composition depth, and $O(Td)$ storage.
\end{exercisechunk}
